\section{绪论}
\subsection{问题定义与应用情境}
应用型研究需要说明真实情境中的问题、利益相关方和约束。应把“研究对象是什么”“希望回答什么问题”“可以取得哪些证据”分开写清楚。选题不能只在题名中添加企业或城市名称而缺少相应资料，也不应预设必须得到正面的结果。
\subsection{文献综述与已有方法}
阅读原始资料，比较方法的假设、适用条件和证据质量，而不是将若干摘要拼接成综述。下面引用真实出版物，仅用于检验中文和英文引文排版\cite{zhou2016,lamport1994}；实际论文须根据选题重新检索和核对。
\subsubsection{概念及分析单位}
为关键概念给出操作性定义，说明分析单位是个体、事件、产品还是组织。样本数、记录数和重复测量次数不应混用。本脚注演示解释性补充\footnote{本文件没有收集企业或个人资料，所有研究场景均为写作提示。}。
\subsection{研究内容与章节安排}
章节安排应服务于问题和证据链。工科研究可围绕需求、设计、实现和测试组织；文经管类可围绕概念、资料、分析和建议组织。学校通知要求各学院按学科特点确定格式，不能把本工作稿的章节清单当成统一要求。
\clearpage\section{研究设计与实施方案}
\subsection{资料来源与使用权限}
报告来源、版本、时间范围和取得方式。企业内部材料应取得适当授权，个人资料应按要求匿名化。无法公开原始材料时，应说明限制，并提供允许披露的变量定义和处理步骤，不能因资料受限而编造数字。
\subsection{技术路线}
图~\ref{fig:route} 展示可修改的技术路线。每一步均应有明确产物，例如研究问题表、资料清单、分析脚本和结果核查记录。
\begin{figure}[htbp]\centering
\begin{tikzpicture}[node distance=6mm,b/.style={draw,minimum width=72mm,minimum height=10mm,font=\small},>=Stealth]
\node[b](a){明确问题、边界与评价准则};
\node[b,below=of a](b){获得资料与必要授权};
\node[b,below=of b](c){实施研究并保留过程记录};
\node[b,below=of c](d){分析证据、复核结果与讨论局限};
\draw[->](a)--(b);\draw[->](b)--(c);\draw[->](c)--(d);
\end{tikzpicture}
\caption{研究组织路线示例}\label{fig:route}
\end{figure}
\subsection{评价指标与数学表达}
指标应与研究目标对应，给出分母、单位和计算规则。例如平均绝对误差可写为：
\begin{equation}\label{eq:mae}\operatorname{MAE}=\frac{1}{n}\sum_{i=1}^{n}|y_i-x_i|.\end{equation}
式中，$x_i$ 为参考量，$y_i$ 为观测量，$n$ 为有效观测数。并非每种应用研究都应使用这一指标，不能为了插入公式而选用不适合的方法。
\subsection{复现与质量控制}
记录软件版本、参数、资料处理规则和异常情况。适用时应使用独立验证、交叉核查或反例分析，区分质量控制与事后挑选结果。研究涉及人类参与者、动物或敏感数据时，须遵守相应伦理与安全要求。
\clearpage\section{结果报告与分析}
\subsection{结果呈现方式}
真实论文应先说明资料规模、缺失和排除，再报告与研究问题直接相关的结果。表~\ref{tab:evidence} 是证据追踪表的结构示例，不包含真实调查结论。
\begin{table}[htbp]\centering\caption{证据追踪表结构示例}\label{tab:evidence}
\begin{tabularx}{\linewidth}{lXX}\toprule 研究问题 & 应保留的证据 & 分析限制\\\midrule
问题一 & 来源记录、原始变量与处理规则 & 样本覆盖与来源偏差\\
问题二 & 方案说明、对照条件与测试记录 & 混杂因素和评价范围\\
问题三 & 分析过程、异常记录与反例 & 外推条件与不确定性\\\bottomrule\end{tabularx}
\end{table}
\subsection{事实、解释与推测}
把“资料显示了什么”“作者如何解释”和“尚需验证的推测”分开。相关性不能直接支持因果结论，个别成功案例也不能证明普遍有效。图表必须有来源与正文引导，避免重复堆放同一组数据。
\subsection{负面结果与异常记录}
如实保留没有达到预期的结果，说明处理方法及其影响。不能把未开展的测试写成已完成，也不能以预期目标充当实际结果。本模板没有实证发现，因此不提供任何性能或政策效果结论。
\clearpage\section{讨论、建议与结论}
\subsection{与已有证据的比较}
讨论应将本研究放回相关文献和应用情境，比较方法及条件差异。若提出实践建议，需要说明实施主体、资源、风险及适用前提，不宜只列空泛口号。
\subsection{局限与未来工作}
说明数据、方法和实施条件带来的限制，区分可在本研究中解决的问题和未来研究计划。若未能完成某项工作，应明确记为未完成，不得使用“验证了”“证明了”等措辞。
\subsection{结论}
结论应回答研究问题并与结果一致，不引入新证据。本文件只演示论文结构和版式，不能作为研究成果提交。正式使用前还须取得所在学院的备案格式，核对检测、答辩和归档要求。
